\label{chap:83-18}

La limite projective \(\Omega_\infty\) conserve l'histoire complète. Pour
obtenir une coordonnée réelle, il faut ajouter une famille compatible
d'intervalles et démontrer trois faits différents : chaque histoire détermine une
unique valeur ; toute valeur du compact déclaré est atteinte ; et les mesures
finies descendent de manière cohérente. Aucun de ces faits ne se déduit
de la seule croissance du nombre d'états. La construction qui suit
démontre conjointement les trois faits pour le système HMT.

La publication archimédienne conserve également une approximation de
l'identité par des espérances conditionnelles cylindriques. La concentration de masse
unitaire sur des préfixes compatibles fournit le passage exact des
cylindres nonadiques à la mesure de Dirac, sans identifier densité,
mesure du support et masse totale.

\begin{HMTScientificOwnerSubordinated}
\input{ampliaciones_sucesoras_20260824/parte_i_ii/owners/03_cilindros_dirac}
\end{HMTScientificOwnerSubordinated}

\section{Système d'intervalles associé aux états finis}

Soit \((S_n,r_{n,m})\) le système projectif de la
\cref{chap:sistema-inverso-continuidad-genealogica}. À chaque
\(a\in S_n\), nous attribuons un intervalle compact non vide
\begin{equation}
 I_n(a)=[\ell_n(a),u_n(a)]\subset\mathbb R.
 \label{p12-83-c18-s1:eq:intervalo-asociado-estado}
\end{equation}
L'affectation est \emph{compatible avec le raffinement} lorsque
\begin{equation}
 r_{n+1,n}(b)=a
 \quad\Longrightarrow\quad
 I_{n+1}(b)\subseteq I_n(a).
 \label{p12-83-c18-s1:eq:anidamiento-intervalos-estados}
\end{equation}
Elle est \emph{uniformément contractante} lorsqu'il existe une suite
\(\delta_n\downarrow0\) telle que
\begin{equation}
 \operatorname{diam}I_n(a)\leq\delta_n
 \qquad(a\in S_n).
 \label{p12-83-c18-s1:eq:contraccion-uniforme-intervalos}
\end{equation}

Pour une histoire \(x=(x_n)_n\), les inclusions
\eqref{p12-83-c18-s1:eq:anidamiento-intervalos-estados} produisent une chaîne
\begin{equation}
 I_0(x_0)\supseteq I_1(x_1)\supseteq I_2(x_2)\supseteq\cdots.
 \label{p12-83-c18-s1:eq:cadena-intervalos-historia}
\end{equation}

\begin{theorem}[Évaluation d'une histoire]
\label{p12-83-c18-s1:thm:evaluacion-unica-historia-u024}
Si l'affectation est compatible et uniformément contractante, alors pour
chaque \(x\in\Omega_\infty\) il existe un unique nombre réel
\begin{equation}
 \operatorname{ev}_{\mathbb R}(x)
 \in\bigcap_{n\geq0}I_n(x_n).
 \label{p12-83-c18-s1:eq:evaluacion-real-historia-u024}
\end{equation}
L'application
\begin{equation}
 \operatorname{ev}_{\mathbb R}:\Omega_\infty\longrightarrow\mathbb R
 \label{p12-83-c18-s1:eq:aplicacion-evaluacion-real-u024}
\end{equation}
est continue.
\end{theorem}

\begin{proof}
Les intervalles de \eqref{p12-83-c18-s1:eq:cadena-intervalos-historia} sont compacts,
non vides et emboîtés. Le théorème des intervalles emboîtés garantit que
leur intersection est non vide. Si \(s,t\) lui appartiennent, alors
\[
 |s-t|\leq\operatorname{diam}I_n(x_n)\leq\delta_n
\]
pour tout \(n\). En faisant \(n\to\infty\), on obtient \(s=t\).

Pour la continuité, soit \(\varepsilon>0\) et choisissons \(n\) avec
\(\delta_n<\varepsilon\). Si \(y\in[x_n]_n\), alors
\(I_n(y_n)=I_n(x_n)\). Les deux évaluations appartiennent à cet intervalle et,
par conséquent,
\[
 |\operatorname{ev}_{\mathbb R}(x)
  -\operatorname{ev}_{\mathbb R}(y)|<\varepsilon.
\]
Le cylindre \([x_n]_n\) est un voisinage ouvert de \(x\).
\end{proof}

La construction n'identifie pas \(x\) à
\(\operatorname{ev}_{\mathbb R}(x)\). Une même coordonnée peut avoir
plusieurs histoires, et une histoire contient davantage de données que n'importe quel point de
la droite.

\section{Couverture cohérente du compact visible}

Soit \(K\subset\mathbb R\) un compact non vide. La construction HMT des
cylindres associe à l'emboîtement et à la contraction les identités de
couverture
\begin{align}
 K&=\bigcup_{a\in S_0}I_0(a),
 \label{p12-83-c18-s1:eq:cobertura-inicial-compacto-u024}\\
 I_n(a)&=
 \bigcup_{\substack{b\in S_{n+1}\\r_{n+1,n}(b)=a}}
 I_{n+1}(b)
 \qquad(a\in S_n).
 \label{p12-83-c18-s1:eq:cobertura-extensiones-compatible-u024}
\end{align}
La seconde égalité est la publication géométrique de la loi de prolongation :
chaque point de \(I_n(a)\) demeure dans un cylindre enfant compatible. Conjointement à
la surjectivité de \(r_{n+1,n}\), déjà établie par la survie des
préfixes, elle produit simultanément des états au-dessus de \(a\) et la
couverture intégrale de \(I_n(a)\).

\begin{theorem}[Surjectivité de l'évaluation]
\label{p12-83-c18-s1:thm:sobreyectividad-evaluacion-compacto}
Les identités construites
\eqref{p12-83-c18-s1:eq:cobertura-inicial-compacto-u024} et
\eqref{p12-83-c18-s1:eq:cobertura-extensiones-compatible-u024} impliquent
\begin{equation}
 \operatorname{ev}_{\mathbb R}(\Omega_\infty)=K.
 \label{p12-83-c18-s1:eq:imagen-evaluacion-compacto}
\end{equation}
\end{theorem}

\begin{proof}
L'inclusion de gauche à droite découle de
\(I_n(x_n)\subseteq I_0(x_0)\subseteq K\).

Fixons \(t\in K\). Définissons
\begin{equation}
 T_n(t):=\{a\in S_n:t\in I_n(a)\}.
 \label{p12-83-c18-s1:eq:estados-que-cubren-punto}
\end{equation}
La couverture initiale rend \(T_0(t)\) non vide. Si \(a\in T_n(t)\),
l'identité \eqref{p12-83-c18-s1:eq:cobertura-extensiones-compatible-u024} fournit au
moins un \(b\in T_{n+1}(t)\) avec \(r_{n+1,n}(b)=a\). Le graphe de ces
extensions est à ramification finie et possède des niveaux non vides à toute
profondeur. Il existe donc une chaîne compatible \(x=(x_n)_n\) avec
\(t\in I_n(x_n)\) pour tout \(n\). L'unicité de l'intersection impose
\(\operatorname{ev}_{\mathbb R}(x)=t\).
\end{proof}

\section{Quotient généalogique archimédien et homéomorphisme avec \texorpdfstring{\(\mathbb R\)}{R}}

Nous définissons la relation
\begin{equation}
 x\sim_{\mathbb R}y
 \quad\Longleftrightarrow\quad
 \operatorname{ev}_{\mathbb R}(x)
 =\operatorname{ev}_{\mathbb R}(y).
 \label{p12-83-c18-s1:eq:relacion-equivalencia-arquimediana}
\end{equation}
Soit
\begin{equation}
 q_{\mathbb R}:\Omega_\infty
 \longrightarrow\Omega_\infty/\!\sim_{\mathbb R}
 \label{p12-83-c18-s1:eq:proyeccion-cociente-arquimediano}
\end{equation}
la projection canonique.

\begin{theorem}[Reconstruction du compact comme quotient]
\label{p12-83-c18-s1:thm:homeomorfismo-cociente-compacto-u024}
Pour le système HMT précédent, il existe un unique homéomorphisme
\begin{equation}
 \overline{\operatorname{ev}}_{\mathbb R}:
 \Omega_\infty/\!\sim_{\mathbb R}
 \xrightarrow{\ \cong\ }K
 \label{p12-83-c18-s1:eq:homeomorfismo-cociente-compacto-u024}
\end{equation}
tel que
\begin{equation}
 \operatorname{ev}_{\mathbb R}
 =\overline{\operatorname{ev}}_{\mathbb R}\circ q_{\mathbb R}.
 \label{p12-83-c18-s1:eq:factorizacion-evaluacion-cociente}
\end{equation}
La fibre au-dessus de \(t\in K\) est
\begin{equation}
 \mathfrak G_t
 :=\operatorname{ev}_{\mathbb R}^{-1}(t),
 \label{p12-83-c18-s1:eq:fibra-genealogias-mismo-valor}
\end{equation}
l'ensemble complet des généalogies qui publient \(t\).
\end{theorem}

\begin{proof}
L'application \(\operatorname{ev}_{\mathbb R}\) est continue, surjective et
constante exactement sur les classes de \(\sim_{\mathbb R}\). Par la
propriété universelle du quotient, elle induit une application continue et bijective
\(\overline{\operatorname{ev}}_{\mathbb R}\). Le domaine est compact, car il est
l'image continue du compact \(\Omega_\infty\), et \(K\) est séparé.
Une bijection continue d'un compact vers un espace séparé est un homéomorphisme.
\end{proof}

Le théorème prouve deux choses et les maintient séparées :
\begin{equation}
 \boxed{
 \begin{array}{c}
 \text{l’espace des histoires est profini et conserve la mémoire,}\\
 \text{son quotient par égalité d’évaluation est le compact visible.}
 \end{array}}
 \label{p12-83-c18-s1:eq:dos-niveles-continuo-u024}
\end{equation}
Il n'est pas affirmé que \(\Omega_\infty\) soit homéomorphe à \(K\).

\section{Atlas généalogique de la droite réelle}

Pour \(m\geq1\), soit \(K_m=[-m,m]\). La même construction, appliquée à chaque
\(K_m\), produit un système \(\Omega_\infty^{(m)}\) qui reconstruit
\(K_m\), avec des inclusions fermées
\begin{equation}
 j_m:\Omega_\infty^{(m)}\hookrightarrow
 \Omega_\infty^{(m+1)}
 \label{p12-83-c18-s1:eq:inclusion-atlas-genealogico}
\end{equation}
qui satisfont
\begin{equation}
 \operatorname{ev}_{m+1}\circ j_m
 =\iota_m\circ\operatorname{ev}_m,
 \label{p12-83-c18-s1:eq:compatibilidad-atlas-real}
\end{equation}
où \(\iota_m:K_m\hookrightarrow K_{m+1}\) est l'inclusion ordinaire.

\begin{proposition}[Exhaustion de la droite]
\label{p12-83-c18-s1:prop:agotamiento-recta-real-u024}
Le quotient archimédien de la colimite stricte
\begin{equation}
 \Omega_\infty^{\mathbb R}
 :=\varinjlim_m\Omega_\infty^{(m)}
 \label{p12-83-c18-s1:eq:colimite-historias-recta}
\end{equation}
est
\begin{equation}
 \varinjlim_m[-m,m]
 =\bigcup_{m\geq1}[-m,m]
 =\mathbb R
 \label{p12-83-c18-s1:eq:colimite-compactos-recta}
\end{equation}
munie de sa topologie usuelle.
\end{proposition}

\begin{proof}
Chaque pièce possède pour quotient \(K_m\) d'après le
\cref{p12-83-c18-s1:thm:homeomorfismo-cociente-compacto-u024} ;
\eqref{p12-83-c18-s1:eq:compatibilidad-atlas-real} entrelace les quotients. La topologie
finale de l'exhaustion par des compacts coïncide avec la topologie usuelle de
\(\mathbb R\) : un sous-ensemble \(U\subseteq\mathbb R\) est ouvert si et seulement
si \(U\cap[-m,m]\) est ouvert dans \([-m,m]\) pour tout \(m\).
\end{proof}

Le passage d'un intervalle borné à la droite n'efface pas les fibres
\(\mathfrak G_t\). Les inclusions \(j_m\) conservent l'histoire et ne font
qu'élargir la carte visible disponible.

\section{Cartes ternaire équilibrée et décimale par blocs}

La représentation ternaire équilibrée d'un entier a la forme
\begin{equation}
 n=\sum_{j=0}^{J}\epsilon_j3^j,
 \qquad
 \epsilon_j\in\{-1,0,+1\},
 \label{p12-83-c18-s1:eq:carta-ternaria-balanceada-u024}
\end{equation}
avec une règle de retenue fixée. Cette carte conserve le signe local et le
quotient de la réduction ; elle ne doit pas être confondue avec le mot non négatif dans
\(\mathbb F_3\).

La carte décimale par blocs utilise
\begin{equation}
 D_N=1000D_{N-1}+u_N,
 \qquad
 u_N\in\{0,\ldots,999\},
 \qquad D_0=0.
 \label{p12-83-c18-s1:eq:concatenacion-base-mil-u024}
\end{equation}
Le bloc \(u_N\), sa classe modulo neuf et la retenue sont des champs distincts.
Puisque \(1000\equiv1\pmod9\),
\begin{equation}
 D_N\equiv\sum_{j=1}^{N}u_j\pmod9.
 \label{p12-83-c18-s1:eq:base-mil-compatible-mod9-u024}
\end{equation}
Cette congruence ne récupère pas l'ordre des blocs ; la séquence
\((u_1,\ldots,u_N)\) demeure dans le registre.

Pour \(M\geq1\), les cellules décimales semi-ouvertes sont
\begin{equation}
 \mathcal D_{M,q}
 :=[q10^{-M},(q+1)10^{-M}),
 \qquad q\in\mathbb Z.
 \label{p12-83-c18-s1:eq:celda-decimal-semiabierta-u024}
\end{equation}
Un intervalle semi-ouvert \([\ell,h)\) est contenu dans une unique cellule si
et seulement si
\begin{equation}
 \left\lfloor10^M\ell\right\rfloor
 =\left\lceil10^Mh\right\rceil-1.
 \label{p12-83-c18-s1:eq:criterio-exacto-celda-decimal-u024}
\end{equation}
La formule avec plafond à l'extrémité droite couvre également le cas où
\(h\) coïncide exactement avec une frontière décimale.

\begin{definition}[Troncature et arrondi]
Pour \(x\geq0\), nous définissons
\begin{align}
 \operatorname{Tr}_M(x)
 &:=10^{-M}\left\lfloor10^Mx\right\rfloor,
 \label{p12-83-c18-s1:eq:truncamiento-decimal-u024}\\
 \operatorname{Rd}_M(x)
 &:=10^{-M}\left\lfloor10^Mx+\frac12\right\rfloor.
 \label{p12-83-c18-s1:eq:redondeo-decimal-u024}
\end{align}
\end{definition}
Ces opérations coïncident sauf lorsque la fraction écartée franchit le
seuil d'une demi-unité à la dernière position. La distinction est sémantique et
arithmétique ; on ne peut substituer l'une à l'autre pour forcer la compatibilité des
préfixes.

\section{Descente des opérations par les fibres d'évaluation}

Soit \(D_\star\subseteq\Omega_\infty\times\Omega_\infty\) le domaine d'une
opération partielle continue
\begin{equation}
 \star:D_\star\longrightarrow\Omega_\infty.
 \label{p12-83-c18-s1:eq:operacion-historias-u024}
\end{equation}
Nous supposons que \(D_\star\) est saturé par les fibres de
\(\operatorname{ev}_{\mathbb R}\times\operatorname{ev}_{\mathbb R}\).

\begin{theorem}[Critère exact de descente]
\label{p12-83-c18-s1:thm:criterio-descenso-operaciones-u024}
Il existe une unique opération visible
\begin{equation}
 \overline\star:
 (\operatorname{ev}_{\mathbb R}\times\operatorname{ev}_{\mathbb R})
 (D_\star)\longrightarrow K
 \label{p12-83-c18-s1:eq:operacion-descendida-visible}
\end{equation}
qui satisfait
\begin{equation}
 \operatorname{ev}_{\mathbb R}(x\star y)
 =\operatorname{ev}_{\mathbb R}(x)
  \ \overline\star\ 
  \operatorname{ev}_{\mathbb R}(y)
 \label{p12-83-c18-s1:eq:entrelazamiento-operacion-descendida}
\end{equation}
si et seulement si
\begin{equation}
 \begin{split}
 &\operatorname{ev}_{\mathbb R}(x)=
  \operatorname{ev}_{\mathbb R}(x'),\\
 &\operatorname{ev}_{\mathbb R}(y)=
  \operatorname{ev}_{\mathbb R}(y')
 \end{split}
 \quad\Longrightarrow\quad
 \operatorname{ev}_{\mathbb R}(x\star y)=
 \operatorname{ev}_{\mathbb R}(x'\star y')
 \label{p12-83-c18-s1:eq:constancia-operacion-sobre-fibras}
\end{equation}
pour tous les couples du domaine.
\end{theorem}

\begin{proof}
Si \(\overline\star\) existe, le membre de droite de
\eqref{p12-83-c18-s1:eq:entrelazamiento-operacion-descendida} ne dépend que des deux
évaluations, et l'on obtient
\eqref{p12-83-c18-s1:eq:constancia-operacion-sobre-fibras}. Réciproquement, cette
constance permet de définir
\[
 s\ \overline\star\ t
 :=\operatorname{ev}_{\mathbb R}(x\star y)
\]
pour tous représentants \(x,y\) d'évaluations \(s,t\). La
condition prouve que la définition ne dépend pas des représentants.
L'unicité est immédiate.
\end{proof}

Les feuillets additif et multiplicatif d'APP fournissent des opérations sur un
support commun, mais leurs retenues et leurs mémoires ne s'identifient pas. Le
théorème précédent permet à la somme et au produit de descendre vers la carte visible
sans exiger que leurs réalisations généalogiques coïncident.

\section{Correspondance entre cardinalité discrète, mesure cylindrique et publication archimédienne}

À chaque profondeur apparaissent trois opérations qui doivent demeurer
séparées :
\begin{align}
 \operatorname{Count}_n(A)&:=|A|,
 \label{p12-83-c18-s1:eq:operacion-conteo-u024}\\
 \operatorname{Measure}_n(A)&:=\mu_n(A),
 \label{p12-83-c18-s1:eq:operacion-medida-u024}\\
 \operatorname{Publish}_n(a)&:=L_n(a).
 \label{p12-83-c18-s1:eq:operacion-publicacion-u024}
\end{align}
Le dénombrement dépend du nombre d'états ; la mesure dépend de poids
compatibles ; la publication dépend du lecteur. Deux fibres de même
cardinal peuvent avoir des poids distincts, et deux états de même publication
peuvent appartenir à des cylindres distincts. L'égalité de l'une de ces trois
sorties n'implique pas l'égalité des deux autres.

\section{Obstructions de l'évaluation archimédienne : emboîtement, continuité, quotient, mesures projectives et descente des opérations}

\begin{criterion}[Réfutation de l'évaluation et de la mesure]
La construction est réfutée si l'un des faits
suivants est vérifié :
\begin{enumerate}
 \item une extension compatible reçoit un intervalle qui n'est pas contenu
       dans l'intervalle de sa restriction ;
 \item les diamètres ne tendent pas vers zéro, mais l'unicité de
       l'intersection est affirmée ;
 \item l'évaluation n'est pas continue dans la topologie des cylindres ;
 \item on affirme que tout point de \(K\) apparaît sans démontrer la couverture
       cohérente \eqref{p12-83-c18-s1:eq:cobertura-extensiones-compatible-u024} ;
 \item l'application induite par le quotient n'est pas bijective ;
 \item l'espace profini est identifié à son quotient visible ;
 \item les inclusions de l'atlas n'entrelacent pas leurs évaluations ;
 \item on utilise \(\lfloor10^Mh\rfloor\) au lieu de
       \(\lceil10^Mh\rceil-1\) pour un intervalle semi-ouvert et l'on perd
       le cas de frontière ;
 \item on confond troncature et arrondi ;
 \item une opération descendue dépend du représentant choisi dans une
       fibre ;
 \item une famille \((\mu_n)_n\) viole
       \eqref{eq:consistencia-proyectiva-medidas-u024} ;
 \item la mesure d'un cylindre dépend du niveau auquel il est raffiné ;
 \item une mesure conditionnelle cesse d'avoir une somme égale à un ou perd sa cohérence ;
 \item la mesure visible \(\nu\) est identifiée aux mesures internes
       \(\mu_t\) ;
 \item on déduit une flèche entre \(\mathbb R\) et
       \(\mathcal E_\infty\) au seul motif que les deux publications partagent
       un antécédent ;
 \item on utilise le dénombrement d'une fibre comme s'il déterminait sa mesure ou sa
       valeur publiée.
\end{enumerate}
\end{criterion}


\section{De l'espace des histoires à la droite réelle}

Soit
\[
 \Psi:\mathbb Z_3^6\xrightarrow{\ \sim\ }
       (\mathbb F_3^6)^{\mathbb N}
\]
la conjugaison de Hensel entre la mémoire initiale et le ruban complet de
blocs de six trits. Concaténer les six coordonnées de chaque bloc donne
un ruban ternaire
\(
 (a_n)_{n\geq 1}\in\{0,1,2\}^{\mathbb N}
\), sur lequel agit l'évaluation
\[
 \operatorname{ev}_3((a_n))
 :=\sum_{n\geq1}\frac{a_n}{3^n}\in[0,1].
\]
Les deux développements d'une extrémité ternaire représentent la même valeur et
demeurent différenciés avant l'application de l'évaluation. Cela ne constitue pas
une ambiguïté de l'objet généalogique : c'est précisément une fibre du lecteur.

On définit
\[
 X_{\mathrm{raw}}:=\mathbb Z\times\mathbb Z_3^6,
 \qquad
 P(m,x):=m+\operatorname{ev}_3(\operatorname{flat}(\Psi x)).
\]
La coordonnée entière fournit l'atlas dénombrable de la droite et la
coordonnée 3-adique conserve l'histoire complète de la partie fractionnaire.

\begin{theorem}[Reconstruction exacte du quotient archimédien]
\label{p12-83-c18-s2:hmtch:thm-cociente-real}
L'application \(P:X_{\mathrm{raw}}\to\mathbb R\) est surjective. Si
\[
 x\sim_P y
 \quad\Longleftrightarrow\quad
 P(x)=P(y),
\]
alors
\[
 \overline P:X_{\mathrm{raw}}/\!\sim_P\;\longrightarrow\mathbb R,
 \qquad [x]\longmapsto P(x),
\]
est une bijection. En conséquence,
\[
 \left|X_{\mathrm{raw}}/\!\sim_P\right|
 =|\mathbb R|
 =2^{\aleph_0}.
\]
\end{theorem}

\begin{proof}
Toute suite ternaire se regroupe de manière unique en blocs consécutifs de
six chiffres. La conjugaison \(\Psi\) attribue à ce ruban une unique mémoire
3-adique initiale. Étant donné \(r\in\mathbb R\), on prend
\(m=\lfloor r\rfloor\) et un développement ternaire de
\(r-m\in[0,1)\) ; la reconstruction précédente produit
\(x\in\mathbb Z_3^6\) avec \(P(m,x)=r\). Ainsi, \(P\) est surjective.

L'application \(\overline P\) est bien définie par la définition de
\(\sim_P\), elle est surjective puisque \(P\) l'est, et injective puisque deux
classes de même image appartiennent à une même fibre. Enfin,
\[
 |\mathbb Z_3^6|
 =\left|(\mathbb F_3^6)^{\mathbb N}\right|
 =729^{\aleph_0}
 =2^{\aleph_0},
\]
et multiplier par la coordonnée dénombrable \(\mathbb Z\) ne modifie pas ce
cardinal. Le calcul coïncide avec celui du quotient, déjà identifié à
\(\mathbb R\).
\end{proof}

\begin{remark}[Clôture survivante]
\label{p12-83-c18-s2:hmtch:rem-version-superviviente}
La complétion brute et le sous-espace survivant
\(X_{\mathrm{surv}}\) sont deux étapes d'une même construction. La loi de
survie produit des couvertures cohérentes à chaque niveau, une contraction
uniforme des diamètres et une prolongation pour toute cellule visible. Le lemme
de la branche infinie produit alors une section au-dessus de chaque point, et
l'évaluation induit, compact par compact, un homéomorphisme du quotient avec
l'intervalle publié. L'exhaustion par \([-m,m]\) reconstruit toute la droite,
en accord avec les
\cref{p12-83-c18-s1:prop:agotamiento-recta-real-u024,p12-83-c18-s2:hmtch:thm-cociente-real}.
\end{remark}

\section{Moyenne archimédienne et double projection du même état}

Soit \(u=(m,\mathbf u)\in
\mathscr O_{\mathbb R,\leftrightarrow}^{\rm HMT}\). Un préfixe
\(\mathbf u_n\) détermine dans la carte fractionnaire un cylindre orienté
semi-ouvert
\[
 I_n(\mathbf u_n)=[a_n,b_n),
\]
avec la convention complémentaire aux extrémités à double développement. Pour
appliquer la compacité, on utilise son adhérence
\(K_n=[a_n,b_n]\). Les compatibilités donnent
\[
 K_{n+1}\subseteq K_n,
 \qquad
 \delta_n:=b_n-a_n\longrightarrow0.
\]
La moyenne archimédienne globale du cylindre est
\[
 \mu_n(u)=m+\frac{a_n+b_n}{2}.
\]
Si \(q\geq n\), alors
\[
 |\mu_q(u)-\mu_n(u)|\leq\delta_n,
\]
de sorte que
\begin{equation}
 \mathfrak R(u):=\lim_{n\to\infty}\mu_n(u)
 \label{p12-83-c18-s2:hmtch:eq-media-arquimediana}
\end{equation}
existe. Ce qui tend vers zéro est le diamètre du cylindre ; la valeur réelle est
la résultante de l'emboîtement, non ce qui s'annule.

\begin{theorem}[Quotient two-way et droite réelle]
\label{p12-83-c18-s2:hmtch:thm-cociente-two-way-real}
Soit
\[
 u\sim_{\mathfrak R}v
 \quad\Longleftrightarrow\quad
 \mathfrak R(u)=\mathfrak R(v).
\]
Sous la couverture cohérente des cartes HMT, l'évaluation induit un
homéomorphisme
\begin{equation}
 \overline{\mathfrak R}:
 \mathscr O_{\mathbb R,\leftrightarrow}^{\rm HMT}/\!
 \sim_{\mathfrak R}
 \xrightarrow{\ \sim\ }\mathbb R.
 \label{p12-83-c18-s2:hmtch:eq-homeomorfismo-two-way-real}
\end{equation}
En particulier, il s'agit d'un morphisme HMT continu : non du résultat du choix
d'une histoire isolée dans chaque fibre.
\end{theorem}

\begin{proof}
La borne des cylindres prouve la continuité de \(\mathfrak R\). La couverture
produit une histoire au-dessus de chaque réel et donc la surjectivité. En passant
au quotient par les fibres, l'application induite est bijective. Dans chaque
carte compacte, une bijection continue vers un intervalle séparé est un
homéomorphisme. La topologie globale du quotient est la topologie finale du
recollement de ces cartes ; la fibre d'évaluation identifie l'extrémité
\(m+1\) d'une carte à l'extrémité \(0\) translatée de la carte suivante.
Le recouvrement résultant est localement fini, de sorte que les
homéomorphismes de carte se recollent en un homéomorphisme global vers
\(\mathbb R\).
\end{proof}
